% titel: Blatt 10, A7: Puffer mit Henderson-Hasselbalch
% skript: Mortimer Kap. 17 (Puffer, Henderson-Hasselbalch)

% Teil I: Werkzeuge, kein Lösungsweg.

Ein Puffer hält den pH nahe konstant und besteht aus einer schwachen Säure
und ihrer konjugierten Base. Die Henderson-Hasselbalch-Gleichung beschreibt
den Zusammenhang.

\begin{werkzeugkasten}[Henderson-Hasselbalch]
\[
  pH = pK_s + \lg\frac{[\text{Base}]}{[\text{Säure}]}.
\]
\begin{itemize}
  \item $pK_s$ der schwachen Säure (hier H$_2$PO$_4^-$, $pK_s{=}7.21$).
  \item \textbf{Base:} konjugierte Base der Säure (hier HPO$_4^{2-}$).
  \item \textbf{Säure:} H$_2$PO$_4^-$.
\end{itemize}
Durch Ein-/Ausrechnen des Verhältnisses bei bekanntem Stoffmengen der
Base erhält man die benötigte Menge der Säure (bzw.\ umgekehrt).
\textit{Skript:} Mortimer Kap. 17 (Puffergleichung).
\end{werkzeugkasten}

\begin{achtung}[Darauf achten]
\begin{itemize}
  \item Volumen kürzt sich heraus, solange man Molverhältnisse der
        Pufferkomponenten einsetzt.
  \item Erst Stoffmenge der vorhandenen Base ausrechnen ($n = c\,V$),
        dann das Ratio $=10^{pH-pK_s}$ in Stoffmenge der Säure umwandeln.
  \item $pH<pK_s \Rightarrow$ $[\text{Säure}]>[\text{Base}]$ (Ratio$<1$).
\end{itemize}
\end{achtung}
